\begin{proposition}[A.4.9; Proved; a second integral central-space computation]\label{push:A49}
The homology formula of A.4.8 can also be computed
on the central spaces, by their normalization
resolutions and the strata of each recorded star.
More precisely, let $Y'\to Y$ be one ordinary
blow-up in the algorithm of A.2.5, with closed
smooth centre $C$ of complex codimension $c$.
Let $D,D'$ be the reduced central supports and
$C_0=C\cap D$. In the derived category of integral
sheaves on $D$ there is an isomorphism
\[
 R(\pi_D)_*\mathbb Z_{D'}\simeq
 \mathbb Z_D\oplus
 \bigoplus_{a=1}^{c-1}(i_0)_*\mathbb Z_{C_0}[-2a],
 \tag{A.49a}
\]
where $\pi_D:D'\to D$ and $i_0:C_0\hookrightarrow D$
are the induced continuous maps. In particular
\[
 H^j(D';\mathbb Z)\simeq H^j(D;\mathbb Z)
 \oplus\bigoplus_{a=1}^{c-1}H^{j-2a}(C_0;\mathbb Z).
 \tag{A.49b}
\]
The maps in this splitting are integral, with unit
diagonal on the exceptional-fibre basis. Thus it
also determines the primitive boundary maps and
the absence of torsion in the normalization
mapping-cylinder computation along the entire
finite algorithm, not just the Euler characteristic.
\end{proposition}
\begin{proof}
We construct the sheaf maps rather than infer a
splitting from fibre ranks. Let $E=\mathbb P_C(N_{C/Y})$
be the exceptional divisor, let $p:E\to C$, and put
$\xi=c_1(\mathcal O_E(1))$. Pullback from $Y$ gives
the degree-zero summand. For $1\leq a\leq c-1$,
pullback by $p$, multiplication by $\xi^{a-1}$,
and the exceptional Gysin morphism give a derived
map
\[
 i_*\mathbb Z_C[-2a]\longrightarrow
 R\pi_*\mathbb Z_{Y'},\qquad i:C\hookrightarrow Y.
\]
The Gysin map exists integrally for a closed complex
divisor, by its oriented tubular Thom class; neither
the centre nor the ambient space must be compact.
At a point of $C$, proper base change identifies
the stalk of the target with
$H^*(\mathbb P^{c-1};\mathbb Z)$.
Restriction of the exceptional normal line is
$\mathcal O_E(-1)$, so the displayed maps send
the generators to
\[
 1,-\xi,-\xi^2,\ldots,-\xi^{c-1}.
\]
This is an integral basis. Off $C$ the map is
the identity. It is therefore a stalkwise
quasi-isomorphism
\[
 \mathbb Z_Y\oplus
 \bigoplus_{a=1}^{c-1}i_*\mathbb Z_C[-2a]
 \xrightarrow{\ \simeq\ }R\pi_*\mathbb Z_{Y'}.
\]

Now restrict this actual sheaf morphism to $D$.
The central support upstairs is the whole inverse
image $\pi^{-1}D$, including when the centre is
horizontal. Proper base change applies to this
Cartesian square. The inverse-image functor for
integral sheaves is exact, and the inverse image
of $i_*\mathbb Z_C$ is
$(i_0)_*\mathbb Z_{C_0}$. This proves (A.49a).
Taking derived global sections proves (A.49b).
This calculation concerns the stratified map of
central spaces. It does not first calculate the
cohomology of either ambient total neighbourhood.

Here is its explicit relation to the requested
normalization mapping cylinders. Resolve
$\mathbb Z_{D'}$ by the alternating ordered-local-branch
complex: its first term is the constant sheaf on
the normalization, its next term is on the normalized
double-branch space, and so on. Its augmentation is
an integral quasi-isomorphism, because at a point
with $s$ branches it is the augmented simplex
complex on those branches. No division by $s!$ is
used; one uses alternating integral cochains, or
equivalently the orientation local system of each
unordered branch set. Proper pushforward of this
resolution computes the left side of (A.49a).
The same resolution for $D$, together with the
centre complex, computes its right side.

The local matrices can be read directly from the
blow-up charts in A.2.5. If the centre lies on $k$
central branches among its $c$ defining coordinates,
the exceptional fibre is $\mathbb P^{c-1}$. The new
exceptional branch supplies that whole projective
space; the strict transform of old branch $i$
supplies its coordinate hyperplane $H_i$.
Intersections supply the corresponding coordinate
projective subspaces. In each even cohomological
degree the normalization differential is the
alternating restriction matrix with entries
$0,1,-1$. The cover consisting of the whole
projective space and its selected hyperplanes
has the elementary simplex contraction, leaving
one generator $\xi^a$ in each degree $2a$,
$0\leq a<c$. For a horizontal centre, every
central branch meeting its generic central
stratum supplies the whole exceptional fibre;
the same simplex contraction applies to these
copies. Higher boundary strata are treated by
the same coordinate restrictions. The global
Gysin maps just constructed glue the surviving
classes with unit coefficient, so this local
cancellation cannot hide an extension or an
integral index.

Starting from the explicitly computed reference
normalization matrices in A.4.8, (A.49a) therefore
computes the normalization complex along every
recorded blow-up or inverse blow-up. A compact
centre has $C_0=C$. A horizontal toric surface
centre retracts to its rational-chain $C_0$;
a horizontal curve centre has $C_0$ a point.
Their cohomology is free and even. Consequently
the central computation preserves odd cohomology
and torsion, adds one to $b_2$ for every new
boundary ray, and adds one to $b_6$ exactly for
a new central ray. It gives the same
$(b_1,b_3,b_5)=(1,7,0)$,
$b_2=8+k+m$ and $b_6=k+m$ as A.4.8.

Independently, stratification of $\pi_D$ into
$D\setminus C_0$ and $C_0$, with fibres a point
and $\mathbb P^{c-1}$, gives
\[
 \chi(D')-\chi(D)=(c-1)\chi(C_0).
\]
The absolute torus-stratum count of A.4.7 is
$\chi(D)=11k+t_4$. Together with the central
group calculation this gives again
$b_4=9k+t_4-2m-1$. Thus the two computations
agree integrally and in Euler number.

The program \path{code/ab_push_central_pushforward.py}
constructs every local alternating projective-subspace
matrix for $2\leq c\leq4$ and $1\leq k\leq c$,
as well as the possible horizontal cases. Its
Smith calculations check the unit cancellations
in every local dimension occurring in the
fourfolds. The sheaf construction proves their
global compatibility for arbitrary input fans.
\end{proof}
